%!TeX root=../Memory.tex

\chapter{Planificació}
\label{ch:planificacio}

Aquest capítol descriu com s'ha planificat i organitzat el projecte \emph{Demochain} al llarg de la seva execució.
S'hi exposa el cicle de vida adoptat, l'estructura del \emph{backlog} amb la seva nomenclatura (èpiques, històries
d'usuari, tasques i \emph{spikes}), les definicions de \aclu{DoR} i \aclu{DoD}, els \emph{sprints} executats, els camins
crítics identificats i la política de versionat del codi.
L'objectiu és deixar constància d'una manera fidel de com s'ha fet la feina --no només del resultat final-- per poder
justificar les decisions de procés que han condicionat el producte entregat.

%-----------------------------------------------------------------------
\section{Cicle de vida del projecte}
\label{sec:cicle-vida}

L'equip ha treballat amb un \textbf{model iteratiu i incremental} inspirat en Scrum~\cite{schwaber2020scrum}, però
adaptat a un context acadèmic amb tres restriccions importants: l'equip és reduït (quatre persones, sense rol de
\emph{product owner} extern), la disponibilitat és parcial (els membres compaginen el projecte amb la resta
d'assignatures del grau) i l'horitzó temporal és tancat (l'entrega final té data fixa).

La conseqüència principal d'aquestes restriccions és que el projecte s'ha hagut de partir en dues fases conceptualment
diferents:

\begin{enumerate}
    \item \textbf{Fase de recerca i fonaments} (febrer-abril de 2026).
    L'equip ha hagut d'adquirir des de zero el coneixement tècnic sobre \emph{blockchain}, Algorand, Ethereum, mètodes
    de votació ordinal i criptografia aplicada.
    Sense aquesta inversió inicial, qualsevol decisió de disseny posterior seria una conjectura.

    \item \textbf{Fase de construcció iterativa} (abril-maig de 2026).
    Amb el coneixement adquirit, s'ha pogut planificar el \emph{backlog} amb èpiques realistes, executar \emph{sprints}
    curts i prendre decisions arquitectòniques justificades.
\end{enumerate}

Aquesta divisió no és estricta: durant la fase de construcció s'han mantingut \emph{spikes} puntuals de recerca (com la
decisió de llenguatge per als \emph{smart contracts}, recollida a l'\ac{ADR} 000) sempre que ha calgut explorar un espai
de solucions nou.

\subsection{Influència de la recerca sobre la implementació}
\label{subsec:influencia-recerca}

La fase de recerca ha condicionat el projecte de manera molt directa.
Les conclusions documentades a~\cite{demochainrepo}, dins la carpeta \texttt{docs/research/}, han portat a les decisions
següents:

\begin{itemize}
    \item L'estudi comparatiu de plataformes \emph{blockchain} (\texttt{docs/research/} i \ac{ADR} 001) va concloure que
    Algorand era la millor opció per a un \ac{MVP} de votació, per la combinació de finalitat ràpida, cost negligible i
    suport per a \emph{smart contracts} expressius.
    Aquesta decisió va condicionar totes les fases posteriors.

    \item L'estudi del mètode de Schulze (\texttt{docs/research/schulze.md}) i l'estudi comparatiu amb altres mètodes de
    votació ordinal va portar a triar Schulze pel seu compliment de propietats formals com la monotonicitat, la
    independència de clons i la consistència amb Condorcet~\cite{schulze2011newmethod}.
    L'algorisme s'implementa amb Floyd-Warshall~\cite{floyd1962algorithm} sobre la matriu de comparació per parelles.

    \item L'estudi dels llenguatges de \emph{smart contracts} d'Algorand (\ac{ADR} 000) va portar a adoptar
    \emph{algopy} amb compilador Puya~\cite{algorandfoundation2024puya} en lloc de \emph{TEAL} o \emph{PyTEAL}, pel seu
    sistema de tipus fort i la integració nativa amb AlgoKit~\cite{algorandfoundation2024algokit}.

    \item L'exploració inicial dels mecanismes de privacitat criptogràfica (xifratge homomòrfic ElGamal, proves
    \acsp{zkSNARK} Groth16, cerimònies de desxifrat amb \emph{trustees}) va revelar una complexitat que excedia el
    pressupost temporal del projecte.
    Aquesta constatació --recollida explícitament al final del primer mes-- ha portat a reduir l'abast de l'\ac{MVP} a
    votació en clar amb verificació pública (vegeu secció~\ref{sec:reduccio-abast}).
\end{itemize}

%-----------------------------------------------------------------------
\section{Organització del backlog}
\label{sec:backlog}

Tot el treball s'ha gestionat amb \textbf{GitHub Projects} sobre el repositori \texttt{demochain}~\cite{demochainrepo}.
Es tracta d'una elecció pragmàtica: GitHub Projects ofereix integració directa amb les \emph{issues}, les \emph{pull
requests} i les revisions de codi, sense introduir cap eina externa.

\subsection{Jerarquia de \emph{issues}}
\label{subsec:jerarquia}

El \emph{backlog} s'estructura en cinc nivells, amb plantilles d'\emph{issue} i etiquetes (\emph{labels}) que en
formalitzen el tipus.
Les plantilles són al directori \texttt{.github/ISSUE\_TEMPLATE/} del repositori i una acció de GitHub
(\texttt{.github/workflows/issue-title.yml}) en valida el títol.

\begin{itemize}
    \item \textbf{Èpica} (\texttt{type/epic}).
    Conjunt de treball que abasta diversos \emph{sprints}.
    Cada èpica té un identificador (\emph{p.
    ex.} E2, E7, E11), un objectiu, una justificació, un abast inclòs i exclòs i un criteri de tancament.

    \item \textbf{Història d'usuari} (\texttt{type/user-story}).
    Funcionalitat de valor visible per a un usuari, formulada en primera persona.
    Cada \ac{US} té identificador format pel codi de l'èpica (\emph{p.
    ex.} E2-\ac{US}1) i un conjunt de criteris d'acceptació en format \emph{Donat/Quan/Aleshores}.

    \item \textbf{\emph{Spike}} (\texttt{type/spike}).
    Investigació limitada en el temps amb un objectiu d'aprenentatge concret i una pregunta a respondre.
    Els resultats d'un \emph{spike} es materialitzen en un document de recerca a \texttt{docs/research/} o en un
    \ac{ADR}.

    \item \textbf{Tasca} (\texttt{type/task}).
    Unitat concreta de treball tècnic, idealment menor a un dia.
    Identificador format pel codi de la \ac{US} mare (\emph{p.
    ex.} E2-\ac{US}1-T3).
    Inclou descripció, criteris d'acceptació i una llista explícita de dependències.

    \item \textbf{Sub-tasca} (\texttt{type/subtask}).
    Subdivisió d'una tasca quan la implementació demana una granularitat més fina.
\end{itemize}

Addicionalment, totes les \emph{issues} es classifiquen segons la seva àrea funcional mitjançant etiquetes
\texttt{area/*}: \texttt{area/smart-contract}, \texttt{area/\textit{front-end}}, \texttt{area/network} (anchoring),
\texttt{area/census}, \texttt{area/wallet}, \texttt{area/test}, \texttt{area/research}, \texttt{area/docs} i
\texttt{area/devops}.

\subsection{Definition of Ready (\acs{DoR})}
\label{subsec:dor}

Una tasca s'ha considerat \emph{ready} per ser introduïda en un \emph{sprint} quan complia simultàniament:

\begin{itemize}
    \item Té un identificador i un títol que segueix el patró \texttt{[TYPE] descripció}.

    \item Té una descripció clara de què cal fer.

    \item Té criteris d'acceptació verificables en format \emph{Donat/Quan/Aleshores}.

    \item Té llistades les dependències --tasques o decisions de les quals depèn-- i totes elles estan resoltes o
    aturades a l'espera d'aquesta tasca.

    \item Està vinculada a una història d'usuari pare i, per transitivitat, a una èpica.

    \item Té assignat almenys un responsable.
\end{itemize}

\subsection{Definition of Done (\acs{DoD})}
\label{subsec:dod}

Una tasca s'ha considerat com a \emph{done} quan complia totes les condicions següents:

\begin{itemize}
    \item Tots els \emph{checkboxes} dels criteris d'acceptació estan marcats com a complerts.

    \item El codi viu en una branca específica i una \ac{PR} està oberta contra \texttt{main}.

    \item La \ac{PR} té un títol que compleix \emph{Conventional Commits}~\cite{conventionalcommits} i és validada per
    la \emph{check} \texttt{PR checks / Conventional Commits title}.

    \item La \ac{PR} passa la \emph{check} \texttt{PR checks / Ruff lint}, que executa \texttt{ruff check} i
    \texttt{ruff format --check} sobre el codi Python del repositori.
    El \emph{linting} del codi TypeScript del client (oxlint) s'executa als \emph{hooks} locals abans del \emph{push}.

    \item La \ac{PR} té com a mínim una revisió aprovatòria d'una persona diferent del darrer \emph{committer}.

    \item Els tests existents continuen passant i s'han afegit tests nous quan la tasca ho requereix.

    \item S'ha fet \textit{merge} de la \ac{PR} a \texttt{main} amb \emph{squash and merge} (vegeu
    secció~\ref{sec:versionat}).
\end{itemize}

%-----------------------------------------------------------------------
\section{Inventari d'èpiques}
\label{sec:epiques}

El \emph{backlog} consta de nou èpiques.
La taula~\ref{tab:epiques} en mostra una visió global, amb el seu estat actual al repositori i la indicació de si la
seva funcionalitat ha quedat coberta per l'\ac{MVP} entregat.

\begin{table}[ht]
\centering
\small
\begin{tabular}{@{}llp{5cm}cc@{}}
\toprule
\textbf{ID} & \textbf{Codi} & \textbf{Títol} & \textbf{Estat} & \textbf{Al MVP} \\
\midrule
\#3   & E1  & Recerca i fonaments                       & Tancada & Sí \\
\#5   & E2  & Propostes i aprovació                      & Tancada & Sí \\
\#6   & E3  & Votació preferencial de l'elecció          & Tancada & Parcial \\
\#7   & E4  & Resultats i verificació del recompte       & Tancada & Parcial \\
\#340 & E5  & Elegibilitat de votants (\acs{ZK})         & Oberta  & No \\
\#10  & E6  & Autenticació del votant i gestió de claus  & Tancada & Parcial \\
\#394 & E7  & Gestió d'organitzacions                    & Tancada & Sí \\
\#413 & E7$^{\star}$ & Ancoratge de resultats a Ethereum & Tancada & Sí \\
\#402 & E11 & Setup del projecte i del repositori        & Oberta  & Sí \\
\bottomrule
\end{tabular}
\caption[Inventari d'èpiques]{Inventari d'èpiques del projecte($^{\star}$).
Hi ha dues èpiques codificades com a E7 al repositori, conseqüència d'una recodificació posterior del \emph{backlog}.}
\label{tab:epiques}
\end{table}

Les marques \emph{Parcial} corresponen a èpiques que es van tancar amb la part bàsica implementada (votació en clar,
recompte amb Schulze, connexió de \emph{wallet}) però que originalment incloïen funcionalitats criptogràfiques que es
van moure a treball futur.
La justificació es detalla a la secció~\ref{sec:reduccio-abast}.

%-----------------------------------------------------------------------
\section{Sprints executats}
\label{sec:sprints}

El projecte s'ha organitzat en iteracions de dues setmanes anomenades \emph{sprints}.
Els tres últims \emph{sprints} estan formalment registrats com a \emph{milestones} al repositori; els tres primers,
executats abans de la consolidació del repositori actual el 17 d'abril de 2026, no hi consten explícitament però sí en
la documentació de seguiment.

\begin{description}
    \item[\textbf{Sprint 1} (febrer-març).] Recerca preliminar.
    Estudi d'Algorand, Ethereum, AlgoKit i mètodes de votació ordinal.
    Setup del primer repositori i prova de concepte sobre LocalNet d'Algorand.

    \item[\textbf{Sprint 2} (març).] Decisions arquitectòniques.
    Diagrames C4 inicials~\cite{brown2018c4}, primer esborrany del contracte \emph{algopy} i selecció definitiva del
    mètode de votació.

    \item[\textbf{Sprint 3} (març-12 abril).] Primer prototip de votació.
    Lògica de propostes i de votació per pluralitat al contracte, mockup de les pantalles principals, configuració de
    \texttt{docker-compose}.
    Entrega intermèdia a la convocatòria del 12 d'abril.

    \item[\textbf{Sprint 4 -- Schulze, Results, Integration} (17--30 abril).] Implementació del mòdul de Schulze en
    TypeScript sobre la matriu desxifrada (issue \#132 i \#133), pàgina de resultats amb guanyador i rànquing (\#136),
    matrius de comparació per parelles col·lapsables (\#139), targeta «El meu vot» al detall de la proposta (\#144) i
    panell de verificació del vot (\#145).
    Aquest \emph{sprint} consolida la presència del repositori actual i incorpora la majoria del codi del client.

    \item[\textbf{Sprint 5 -- Testing, Refinament i Robustesa} (1--14 maig).] Cobertura de proves unitàries del domini
    de propostes (\#458), tests dels \emph{mappers} de \emph{queries} Algorand (\#457), tests d'integració del
    \emph{governance client} contra LocalNet (\#456), correcció del \emph{bug} d'accés a transaccions via Lora (\#463) i
    refinament de la pantalla de votació amb edició d'hora/minuts amb el teclat (\#470).

    \item[\textbf{Sprint 6 -- Documentació, Demo i Entrega} (15--25 maig).] Redacció de la memòria, preparació de la
    presentació oral, retrospectiva i empaquetament final amb la versió \texttt{v1.0.0}.
\end{description}

Els \emph{sprints} 4 a 6 consten com a \emph{milestones} actius al repositori, amb el detall complet de les
\emph{issues} associades a cada un.

%-----------------------------------------------------------------------
\section{Camins crítics}
\label{sec:camins-critics}

L'arquitectura del sistema imposa unes dependències estrictes entre components que han condicionat l'ordre de la
implementació.
Identificar-les amb antelació ha estat clau per evitar bloquejos i, en alguns casos, per prioritzar feina aparentment
menor que altres tasques necessitaven.

\subsection{Camí crític 1: Recerca \textrightarrow{} contracte Algorand}
\label{subsec:cami-critic-1:-recerca-contracte-algorand}

Cap línia de codi del contracte podia escriure's abans d'entendre el model d'execució de l'\ac{AVM}, l'estructura de
\emph{boxes} d'emmagatzematge, el sistema de tipus d'\emph{algopy} i les convencions
\ac{ARC}-4~\cite{algorandfoundation2023arc4}.
Aquesta dependència és la justificació de la \emph{fase de recerca} prèvia descrita a la secció~\ref{sec:cicle-vida}.
La duració d'aquesta fase ha tingut un cost d'oportunitat alt --un mes sense codi entregat-- però sense aquest
coneixement el contracte hauria estat ple d'errors de model.

\subsection{Camí crític 2: Contracte Algorand \textrightarrow{} client web}
\label{subsec:cami-critic-2:-contracte-algorand-client-web}

El client només pot interactuar amb el contracte un cop es genera el fitxer de descripció
\ac{ARC}-56~\cite{algorandfoundation2023arc56} a partir del contracte compilat.
Aquest fitxer, en el repositori \texttt{client/src/algorand/Demochain.arc56.json}, conté l'\ac{ABI} dels mètodes públics
i és consumit pel client TypeScript per construir i signar transaccions amb \emph{algosdk}.

La conseqüència pràctica és que el client no podia avançar més enllà del \emph{mockup} estàtic fins que el contracte
tenia un primer conjunt de mètodes estables.
La planificació ha donat prioritat, dins de cada \emph{sprint}, a tancar primer els mètodes del contracte de l'àrea
funcional corresponent abans que les pantalles del client que els consumeixen.
Quan, ocasionalment, el contracte introduïa canvis en una \ac{PR}, el client havia de regenerar l'\ac{ABI} i adaptar els
\emph{mappers}.

\subsection{Camí crític 3: Contracte Algorand \textrightarrow{} ancoratge}
\label{subsec:cami-critic-3:-contracte-algorand-ancoratge}

El servei d'ancoratge llegeix directament les \emph{boxes} del contracte d'Algorand per construir l'estat canònic de
l'elecció.
Concretament, llegeix les \emph{boxes} amb prefix \texttt{pr\_} (propostes) i \texttt{eb\_} (paperetes d'elecció), tal
com s'implementen al contracte.
Si l'esquema d'aquestes \emph{boxes} canvia, el servei d'ancoratge no pot interpretar les dades.

L'ancoratge, per tant, no es podia desenvolupar de manera realista fins que el contracte tenia un esquema
d'emmagatzematge estabilitzat.
La planificació ha tractat aquesta dependència posant el codi de l'ancoratge en una èpica pròpia (\#413) executada un
cop el contracte ja tenia versió estable, i ha inclòs a la \ac{US} E7-\ac{US}1 (\#414) la tasca explícita
d'\emph{adaptar} la capa d'ancoratge importada d'un repositori extern al nou contracte.

\subsection{Camí crític 4: Infraestructura local de desenvolupament}
\label{subsec:cami-critic-4:-infraestructura-local-de-desenvolupament}

Tots tres components anteriors necessiten un entorn local funcional.
Aixecar un node Algorand LocalNet amb \emph{indexer}, \emph{conduit} i base de dades PostgreSQL no és una tasca trivial.
El fitxer \texttt{docker-compose.yml} del repositori integra tots els serveis i automatitza la compilació i desplegament
del contracte i el seu registre a \texttt{client/.env.local}.

Aquesta infraestructura ha estat un camí crític silenciós: cada vegada que un membre de l'equip incorporava una nova
versió del codi havia de poder aixecar tota la pila en pocs minuts.
La inversió en automatització d'aquest desplegament ha estat un dels factors que ha permès una velocitat sostinguda
durant la fase de construcció.

%-----------------------------------------------------------------------
\section{Paquets de treball}
\label{sec:paquets-treball}

Els paquets de treball es deriven de l'estructura del repositori i de les àrees funcionals etiquetades a les
\emph{issues}.
Cada paquet té un responsable principal --el mateix membre de l'equip que ha liderat l'àrea durant tot el projecte--
però totes les \acsp{PR} han passat per revisió creuada.

\begin{table}[ht]
\centering
\small
\begin{tabular}{@{}p{3cm}p{4.2cm}p{4.5cm}@{}}
\toprule
\textbf{Paquet} & \textbf{Ubicació al repo} & \textbf{Responsable principal} \\
\midrule
Smart contract Algorand    & \texttt{voting-contract/}     & Marc Link Cladera \\
Client web                 & \texttt{client/}              & Jordi Vanyó Mestre \\
Ancoratge a Ethereum       & \texttt{network/anchoring/}   & Dylan Luigi Canning \\
Smart contract notari      & \texttt{network/ethereum/}    & Dylan Luigi Canning \\
Infraestructura local      & \texttt{docker-compose.yml}, \texttt{conduit.yml}, \texttt{Makefile} & Jordi Vanyó / Dylan Luigi \\
Documentació i \emph{backlog} & \texttt{docs/}, GitHub Projects & Antonio Contestí Coll \\
\bottomrule
\end{tabular}
\caption{Distribució de paquets de treball.}
\label{tab:paquets}
\end{table}

%-----------------------------------------------------------------------
\section{Política de versionat del codi}
\label{sec:versionat}

S'ha adoptat una política de \textbf{\emph{feature branches} + \ac{PR} + \emph{squash and merge}}, complementada amb
\emph{Conventional Commits}~\cite{conventionalcommits} i \emph{git hooks} locals.
Les regles concretes estan documentades a \texttt{CONTRIBUTING.md} del repositori i es resumeixen a continuació.

\subsection{Estructura de branques}
\label{subsec:estructura-de-branques}

\begin{itemize}
    \item \texttt{main}.
    Branca principal.
    Conté codi sempre desplegable.
    Cap \emph{push} directe està permès; tots els canvis hi arriben via \ac{PR}.

    \item \texttt{dev}.
    Branca d'integració per a treball en curs i artefactes de documentació que encara no estan a punt per ser publicats
    a \texttt{main} (com els capítols de la memòria).

    \item \texttt{feat/<num>}, \texttt{fix/<num>}, \texttt{docs/<num>}, etc. Branques de característica, una per cada
    \emph{issue}.
    El número correspon a l'identificador GitHub de l'\emph{issue}.
    \emph{P.
    ex.} la branca \texttt{feat/392} resol l'\emph{issue} \#392.
\end{itemize}

\subsection{Convencions de \emph{commit}}
\label{subsec:convencions-de-commit}

Els missatges de \emph{commit} segueixen l'estàndard \emph{Conventional Commits}~\cite{conventionalcommits} amb el
format:

\begin{quote}
\texttt{<type>(<scope>): descripció}
\end{quote}

Els tipus permesos són: \texttt{feat}, \texttt{fix}, \texttt{docs}, \texttt{style}, \texttt{refactor}, \texttt{test},
\texttt{chore}, \texttt{ci}, \texttt{build}, \texttt{perf} i \texttt{revert}.
El compliment es valida en dos punts:

\begin{itemize}
    \item \emph{Localment}, mitjançant els \emph{hooks} de Git (\texttt{.githooks/commit-msg},
    \texttt{.githooks/pre-commit}, \texttt{.githooks/pre-push}) que es registren executant \texttt{make hooks}.

    \item \emph{Remotament}, mitjançant l'acció de GitHub \texttt{.github/workflows/pr-checks.yml}, que comprova el
    títol de cada \ac{PR}.
\end{itemize}

\subsection{Procés de \emph{Pull Request}}
\label{subsec:proces-de-pull-request}

\begin{enumerate}
    \item Es crea una branca \texttt{feat/<num>} a partir de \texttt{main}.

    \item S'hi fan els \emph{commits} que calguin.

    \item En fer \emph{push}, els \emph{hooks} locals executen Ruff (codi Python) i Biome (codi TypeScript) per
    assegurar el format.

    \item S'obre una \ac{PR} contra \texttt{main}.
    La plantilla de \ac{PR} enllaça automàticament l'\emph{issue} corresponent mitjançant la fórmula \texttt{Closes
    \#<num>}.

    \item GitHub Actions executa les \emph{checks} de format i nomenclatura.

    \item Almenys un altre membre de l'equip revisa la \ac{PR} i n'aprova el contingut.

    \item Es fa \emph{squash and merge} per mantenir un historial lineal i una correspondència 1:1 entre \emph{commits}
    a \texttt{main} i \emph{issues} tancades.
\end{enumerate}

Una conseqüència d'aquesta política, observada durant el projecte, és que ningú pot fer \emph{merge} de la seva pròpia
\ac{PR}: la regla de \emph{branch protection} de GitHub exigeix una aprovació d'algú diferent de l'últim \emph{pusher}.
Aquesta restricció ha estat un mecanisme natural de revisió creuada i ha contribuït a la qualitat del codi entregat.

%-----------------------------------------------------------------------
\section{Reducció d'abast: de l'MVP criptogràfic a l'MVP en clar}
\label{sec:reduccio-abast}

Una de les decisions de planificació més rellevants del projecte ha estat la \textbf{reducció d'abast} acordada al final
del primer mes de desenvolupament.

L'estat objectiu inicial, recollit a les èpiques E3, E4, E5 i E6 del \emph{backlog}, era un sistema amb:

\begin{itemize}
    \item Vots xifrats amb ElGamal exponencial homomòrfic.

    \item Proves \acsp{zkSNARK} Groth16 per acreditar l'elegibilitat del votant sense revelar-ne la identitat.

    \item Cerimònia de generació distribuïda de claus (DKG) i de desxifrat amb \emph{trustees} i proves Chaum-Pedersen.

    \item Integració amb un \ac{IdP} extern (Keycloak) per a l'autenticació prèvia del votant.
\end{itemize}

Aquesta arquitectura estava documentada com a estat objectiu i les èpiques corresponents s'havien redactat amb el detall
necessari (vegeu les descripcions de \#10, \#340, \#7 i \#6 al repositori).

Després d'una sessió d'estimació realista cap al final del Sprint~3, l'equip va concloure que la criptografia avançada
\textbf{no era assolible} dins l'horitzó del projecte sense comprometre la resta de l'entrega.
La decisió adoptada va ser:

\begin{enumerate}
    \item Reduir l'\ac{MVP} a votació \textbf{en clar} amb verificació pública via consulta directa a la cadena Algorand
    i ancoratge \textit{K-of-N} a Ethereum.

    \item Tancar les èpiques E3, E4 i E6 amb el subconjunt de funcionalitats efectivament implementades --votació
    preferencial, recompte Schulze, connexió de \emph{wallet}-- i deixar la resta d'\acsp{US} obertes com a treball
    futur explícit.

    \item Mantenir oberta l'èpica E5 sencera, ja que tota la seva funcionalitat depèn de la capa criptogràfica.
\end{enumerate}

Aquesta decisió ha estat un dels re-aprenentatges més importants del projecte.
Vista en retrospectiva, és coherent amb les recomanacions de la literatura sobre gestió de projectes: davant un objectiu
ambiciós i un horitzó tancat, és preferible entregar un subconjunt funcional i sòlid que un superconjunt incomplet o
inestable.
La criptografia avançada queda identificada explícitament al capítol de conclusions com a línia de treball futura amb el
detall necessari perquè un altre equip pugui continuar-hi.
